\documentclass{article}

\usepackage{fancyhdr} % Required for custom headers
\usepackage{lastpage} % Required to determine the last page for the footer
\usepackage{extramarks} % Required for headers and footers
\usepackage[usenames,dvipsnames]{color} % Required for custom colors
\usepackage{graphicx} % Required to insert images
\usepackage{listings} % Required for insertion of code
\usepackage{courier} % Required for the courier font
\usepackage{multirow}
\usepackage{hyperref}
\usepackage{amsmath}
\usepackage{amssymb}
\usepackage[utf8]{inputenc}


% Margins
\topmargin=-0.45in
\evensidemargin=0in
\oddsidemargin=0in
\textwidth=6.5in
\textheight=9.0in
\headsep=0.25in

\linespread{1.1} % Line spacing

%----------------------------------------------------------------------------------------
%	CODE INCLUSION CONFIGURATION
%----------------------------------------------------------------------------------------

\definecolor{MyDarkGreen}{rgb}{0.0,0.4,0.0} % This is the color used for comments
\lstloadlanguages{c} % Load Perl syntax for listings, for a list of other languages supported see: ftp://ftp.tex.ac.uk/tex-archive/macros/latex/contrib/listings/listings.pdf
\lstset{language=[sharp]c, % Use Perl in this example
        frame=single, % Single frame around code
        basicstyle=\small\ttfamily, % Use small true type font
        keywordstyle=[1]\color{Blue}\bf, % Perl functions bold and blue
        keywordstyle=[2]\color{Purple}, % Perl function arguments purple
        keywordstyle=[3]\color{Blue}\underbar, % Custom functions underlined and blue
        identifierstyle=, % Nothing special about identifiers                                         
        commentstyle=\usefont{T1}{pcr}{m}{sl}\color{MyDarkGreen}\small, % Comments small dark green courier font
        stringstyle=\color{Purple}, % Strings are purple
        showstringspaces=false, % Don't put marks in string spaces
        tabsize=5, % 5 spaces per tab
        %
        % Put standard Perl functions not included in the default language here
        morekeywords={rand},
        %
        % Put Perl function parameters here
        morekeywords=[2]{on, off, interp},
        %
        % Put user defined functions here
        morekeywords=[3]{test},
       	%
        morecomment=[l][\color{Blue}]{...}, % Line continuation (...) like blue comment
        numbers=left, % Line numbers on left
        firstnumber=1, % Line numbers start with line 1
        numberstyle=\tiny\color{Blue}, % Line numbers are blue and small
        stepnumber=5 % Line numbers go in steps of 5
}

\newcommand{\horrule}[1]{\rule{\linewidth}{#1}}

% Creates a new command to include a perl script, the first parameter is the filename of the script (without .pl), the second parameter is the caption
\newcommand{\perlscript}[2]{
\begin{itemize}
\item[]\lstinputlisting[caption=#2,label=#1]{#1.cs}
\end{itemize}
}

\begin{document}

\begin{tabular}{l l}
\multirow{5}{*}{\includegraphics[width=2cm]{escudo.png}} & Universidad del Istmo de Guatemala \\
 & Facultad de Ingenieria \\
 & Ing. en Sistemas \\
 & Informatica 1 \\
 & Prof. Ernesto Rodriguez - \href{mailto:erodriguez@unis.edu.gt}{erodriguez@unis.edu.gt} \\
\end{tabular}
\\

\begin{center}
        \horrule{0.5pt}
        \huge{Hoja de trabajo \#2} \\
        \large{Samantha Rodas Chaluleu}\\
        \large{Adalí Garrán Jiménez}\\
        \large{6 de agosto, 2019} \\
        \horrule{1pt}
\end{center}
\section*{Ejercicio \#1}
Demostrar utilizando inducci\'on:
\[
        \forall\ n.\ n^3\geq n^2
\]
\\donde $n\in\mathbb{N}$\\
\\
\large{Caso Base:}\\
n=0
\[
0^3\geq 0^2
\]
\[
0\geq 0
\]\\
\large{Caso Inductivo:}\\
n=n+1\\
\[
n^3\geq n^2 
\] \begin{center}(Hipótesis Inductiva)\end{center}
\[
(n+1)^3\geq (n+1)^2
\]
\[
(n+1)(n+1)(n+1)\geq (n+1)(n+1)
\]
\[
(n+1)\geq \frac{(n+1)(n+1)}{(n+1)(n+1)}
\]
\[
n+1\geq 1
\]
\[
n\geq 1-1
\]
\[
n\geq 0
\]

\section*{Ejercicio \#2}
Demostrar utilizando inducci\'on la desigualdad de Bernoulli:
\[
        \forall\ n.\ (1+x)^n\geq nx
\]
\\donde $n\in \mathbb{N}$, $x\in \mathbb{Q}$ y $x\geq -1$\\
\\
\large{Caso Base:}\\
n=0
\[
        (1+x)^0\geq 0x
\]
\[
        1\geq 0
\]\\
\large{Caso Inductivo:}\\
n=n+1\\
\[
        (1+x)^{n}\geq nx
\]\begin{center}(Hipótesis Inductiva)\end{center}
\[
        (1+x)^{n+1}\geq (n+1)x
\]
\[
        (1+x)^n*(1+x)\geq (n+1)x
\]
\[
        (1+x)^n*(1+x)\geq nx+x
\]
\[
        (1+x)^n+x(1+x)^n\geq nx+x
\]
\[
        x(1+x)^n+[(1+x)^n\geq nx]+x
\]\begin{center}Asumiendo que $[(1+x)^n\geq nx]$ es verdadero (Hipótesis Inductiva) que:\end{center}
\[
        x(1+x)^n\geq x
\]
\[
        (1+x)^n\geq \frac{x}{x}
\]
\[
        (1+x)^n\geq 1
\]
\[
        log(1+x)^n\geq log(1)
\]
\[
        n[log(1)+log(x)]\geq log(1)
\]
\[
        n[0+log(x)]\geq 0
\]
\[
        n\geq \frac{0}{log(x)}
\]
\[
        n\geq 0
\]
\end{document}
